\section{Rescale the local tests and remove the intermediate weight}

The starting source is the frozen sixth-round construction, with builder
SHA256
\texttt{3f63cb782d52857446cf39f6e545aeb74f380d1347826eafa888cbb487738d51}.
The selected integrated source has SHA256
\texttt{78b54bcde24ff8a91091e102a4734c49878699f071d9230e5ce7a70ffd2114d0}.
The change in this section removes one surviving rational coordinate
from each of the three detectors. It leaves the orders $1,12,45$, all
sixteen underlying field tests, the logarithms, height scales, contact
coordinate, and pole predicate unchanged. The global arithmetic
qualifications are inherited; the construction is not asserted to prove
those inputs independently.

\subsection{The normalization freedom in the actual local lemma}

Retain
\[
 t=A\lambda,\qquad\delta=98t-1,\qquad
 F(t)=t(t-1/49)(t-1/50)(t-1/2450)(t-32/1225),
\]
and retain the semantic quantity
\[
 R=\frac{F(t)}{A^5\delta^6}.
\]
The three scales are still $1,C_0^2,(C_0+C_1)^2$. On the source domain
they are positive rational squares and at least one. The positive chart
for $\lambda$ gives $t\ge32/1225$, hence $\delta\ge39/25>0$.

At every prime consumed by the retained five-contact height argument,
$\delta$ is a unit. The values at the five contacts are
\[
 -1,\quad1,\quad24/25,\quad-24/25,\quad39/25.
\]
Their prime support is already contained in
$\mathcal S=\{2,3,5,7,13,17\}$. In particular $13$ remains exceptional;
we do not claim that the guard is a unit there. A pole of $t$ is not a
consumed contact outside this set: its image under
$\phi(t)=-49+1/t$ reduces to $-49$, distinct from the five height contacts.

Let $D_j$ be an existing test of order $m_j$. In the arbitrary-order
one-place lemma, replace its semantic test by
\[
 D'_j=\delta D_j.
\]
No coordinate for this product is required. After the original
normalization, the actual factor is
\[
 1+R\left(\frac{D_j/\delta^{m_j}}{R^{m_j}}\right)^2
   =1+\frac{e^{2m_j-1}}{\delta}D_j^2,
 \qquad e=\frac{(A\delta)^5}{F(t)}.
\]
Thus use the new weights
\[
 q_m=\frac{e^{2m-1}}{\delta}
\]
in place of $e^{2m-1}/\delta^3$. The rational function $e$ is useful
for explaining the identity, but no longer needs its own coordinate.

This is an application of the actual proof of the one-place lemma,
\texttt{manuscripts/quartic\_v10.tex}, lines 11031--11061, in the pinned
upstream source. Its field inputs are arbitrary elements of $K$; it
requires no integrality assumption on $D_j/\delta^m$ before the
conclusion. At a consumed prime the unit property gives
\[
 v_{\mathfrak p}(D_j/\delta^m)=v_{\mathfrak p}(D_j),
\]
so the conclusion is exactly the original order requirement for $D_j$.
All sixteen orders still hold at one common place. No new formal-group
claim, prime-selection theorem, or logarithm precision assertion is
being substituted for the retained one-place argument.

\subsection{Compute the weights without storing $e$}

Keep the existing $U=A\delta$ and its fourth power. Define directly
\[
 q_1F=A U^4,\qquad E=q_1\delta^2.
\]
These give
\[
 q_1=\frac{A^5\delta^4}{F}=\frac e\delta,
 \qquad E=\frac{A^5\delta^6}{F}=R^{-1},
 \qquad q_m=E^{m-1}q_1^m.
\]
The value $\delta^2=9604t^2-196t+1$ is affine in the already stored
$t^2$ and $t$. The implementation requests its square through the
common-subexpression routine, which reuses the existing quadratic
part. It adds no new coordinate for this square.

The order-one product evaluation already requires $q_1^2,q_1^3,q_1^4$.
Using its actual stored cube, compute the two other weights in seven gates:
\begin{align*}
 a&=E^2,& b&=aE=E^3,& c&=bq_1^3,& d&=c^2,\\
 q_{12}&=d^2/E,& h&=q_{12}^2,&q_{45}&=h^2/q_1^3.
\end{align*}
The exponents in $(E,q_1)$ are respectively
\[
 (2,0),(3,0),(3,3),(6,6),(11,12),(22,24),(44,45).
\]
Equivalently, their exponents in the explanatory pair $(e,\delta)$ are
\[
 (2,2),(3,3),(6,0),(12,0),(23,-1),(46,-2),(89,-1).
\]
The last two outputs are therefore exactly the required new weights.
The same grouping of the field factors is retained: one mixed $(1,12)$
block, two four-test order-one blocks, and two three-test order 45 blocks.
All shared symmetric coefficients and the mixed-block identity remain
unchanged.

The old implementation separately stored $e$, $\delta^3$, and the
quotient $q_1=e/\delta^3$. The new one defines $q_1$ directly and never
uses a stored $e$ or $\delta^3$. The latter three coordinates were
already removed by the previous final elimination stage, so their raw
removal does not itself constitute a final saving. The net measured
saving is the other three quotient coordinates, one per detector.

\subsection{Nonzero domains and quadratic degree}

On the positive chart, all factors of $F$ are nonnegative. The relation
$q_1F=AU^4$, whose right side is strictly positive, excludes $F=0$.
Consequently $F,q_1,E$ are positive. Every intermediate in the displayed
chain is positive, and its two quotient denominators $E$ and $q_1^3$
are nonzero. Each auxiliary coordinate is uniquely determined by the
preceding coordinates. Eliminating this chain recovers the exact rational
weight formulas; conversely those formulas uniquely extend every point
of the established nonzero domain.

Every displayed gate equation has degree at most two in stored
coordinates. For example $q_1F=AU^4$ uses the already stored fourth
power, $Eq_{12}=d^2$, and $q_1^3q_{45}=h^2$ uses the already stored cube.
The source checker verifies the actual square/cube aliases reused by the
order-one evaluation, rather than merely counting them as free.

Each factor $1+q_mD_j^2$ is at least one at both real embeddings. If
$W=W_0+W_1\sqrt2$ is their product, then $W_0\ge1$ and $N(W)>0$.
The stored norm quotient and pole input remain
\[
 n=\frac{N(W)}{2W_0},\qquad B=\frac{n^2}{E}
   =\frac{R}{\operatorname{Tr}(W^{-1})^2}.
\]
Their denominators are nonzero. The pole input uses exactly the old
$R$, not a modified rational function differing by an unchecked factor.
The normalization changes the existential detector witnesses; it is
justified by soundness and completeness transfer, not by a claimed
bijection between every old and new pole witness tuple.

\subsection{Completeness at contacts, other primes and infinity}

The preparation and order of choices remain those of the source.
First fix the input, scalar preparation, four root points and rational
square scales. The source real-progression argument supplies bounds on
the unscaled $D_j$ independent of future progression moduli. Its finite
neighborhood preparation supplies common lower valuation bounds on
$D_j$ above $2,3,17$. The exact relevant passage is the pinned
\texttt{quartic\_v10.tex}, lines 11220--11241. The same prime selection
and primary progression $\Pi^{44}\mid h$ are then used. Alternate
points still require only formal depth one.

At a selected contact prime, $R$ has valuation one and $\delta$ is a
unit. Multiplying the old factor's nonconstant term by $\delta^2$
therefore does not change its valuation. At the successful place every
factor is one modulo the prime. At either place every inverse factor
has residue zero or one, because its nonconstant term has odd valuation.
Hence the reciprocal trace has residue one or two and is a unit.
The new $B$ has the same positive odd contact valuations as required.

The exact product formula for $R$ is unchanged:
\[
 R=\frac{\prod_{\theta\in\Theta}(u-(\theta/A)v)}
          {v^5\delta^6},\qquad v=D_0^2w^2.
\]
The denominator is a rational square and the numerator consists of the
selected contact primes, with the existing incidence convention for
repeated forms. At every other prime $v_p(R)$ is even, so $v_p(B)$ is
even regardless of the trace valuation. Thus the reduced denominator
of $B$ is a square. In particular a new zero of $\delta$ away from
contacts is harmless: its contribution to $R$ has even order. The
fixed Gaussian classes still make $R$, and hence $B$, square in
$\mathbb Q_2,\mathbb Q_3,\mathbb Q_{17}$.

For fixed $A$ and $\lambda\to+\infty$,
\[
 e\longrightarrow98^5A^5>0,\qquad
 \delta\sim98A\lambda,\qquad q_m=O(\lambda^{-1}).
\]
The uniform bounds on the unscaled tests therefore give
\[
 W=1+O(\lambda^{-1}),\qquad
 \operatorname{Tr}(W^{-1})=2+O(\lambda^{-1}),\qquad
 B\sim\frac{1}{4\cdot98^6A^6\lambda}.
\]
This is slower decay than before, but it still gives $0<B<17$ for
sufficiently large prescribed $\lambda$. The actual preparation-zero
pole converse requires this bound; no stronger real assertion is used.

At a fixed small prime $p=2,3,17$, put
\[
 \nu=v_p(v),\qquad\eta=v_p(A),\qquad\alpha=v_p(98).
\]
The scale and $\alpha$ are fixed before choosing large $\nu$.
Since $u$ is a selected prime, $v_p(u)=0$. For all sufficiently large
$\nu$, dominance of the leading terms gives
\[
 v_p(\delta)=\eta+\alpha-\nu,\quad
 v_p(e)=5(\eta+\alpha),\quad
 v_p(q_m)=\nu+(10m-6)(\eta+\alpha).
\]
If $v(D_j)\ge-c_p$ on the fixed neighborhood, then
\[
 v(q_mD_j^2)\ge\nu+(10m-6)(\eta+\alpha)-2c_p
             \longrightarrow+\infty.
\]
Thus the trace is as close to two as necessary, including at the
ramified prime two. Since $R$ is unchanged,
\[
 v_p(B)=\nu-6(\eta+\alpha)-2v_p(2).
\]
Prescribed high divisibility of $w$ gives the same required bounds
$v_2(B)\ge3$, $v_3(B)\ge1$, $v_{17}(B)\ge2$. The unchanged
preparation-zero bridge therefore supplies its four pole witnesses.
This completes the changed local and real completeness checks.

\subsection{Actual counts and verification boundary}

The isolated replacement, with all other sixth-round formulas and
$S_0=3^{3072}$ retained, compiles as
\[
 \text{raw }404/372\ \longrightarrow\
 \text{generic }393/361\ \longrightarrow\
 \boxed{\text{final }393/361}.
\]
The pairs count rational coordinates and residuals. No further
elimination is needed after the generic compiler. In particular the
unit-scale fourth-power elimination is now carried out there, and no
old final-elimination saving is counted twice.

With the separately proved scalar outer-constant change and selected
anchor $S_0=3^{2790}$, the integrated source compiles as
\[
 403/371\ \longrightarrow\
 \boxed{392\text{ coordinates and }360\text{ residuals}}.
\]
Again there are zero further eliminations. The fixed-parameter quartic
has 26,301 ordinary-coordinate monomials, degree four and primitive
integer content one. Its coefficient balancing and anchor margins are
separate transformations with their own exact certificates and proofs.

The independent checker \texttt{check\_detector\_round7.py} imports
neither producer nor compiler. On the frozen integrated source it passes
384 source assertions, seven Laurent exponent identities, three semantic
normalization identities, and two expanded mixed-block identities.
It reconstructs the retained $G_{12}/G_2$ logarithms, all test aliases,
all shared coefficients and product evaluations, the direct $q_1$
equation, $E=R^{-1}$, every weight gate, actual reused powers, norm and
pole consumers, the contact and positivity data, and the lifted-degree
bound. The global arithmetic implications remain qualified as in the
source; these finite checks do not prove those inherited results.

\subsection{Rejected searches and the useful obstruction}

A bounded exact exponent search for six-node replacements of the old
weight chain found no candidate in its stated restricted class. This is
not a general minimality result. There is nevertheless a useful simple
obstruction in the old monomial-gate model. With the existing monomials
$e,\delta,\delta^2,\delta^3,E,q_1^2,q_1^3,q_1^4$, where the old
$q_1=e/\delta^3$, the functional $L(a,b)=13a+3b$ on exponent pairs
has all input values in $[0,16]$. A quadratic quotient operation
$A+B-C$ can at most double the width of the current interval.
After six operations its maximum cannot exceed 1024, whereas
$L(q_{41})=1044$. Thus a six-node monomial chain cannot even produce
that high-order weight from these inputs. The search also considered
lower orders and shared outputs, without claiming an exhaustive result
for arbitrary rational circuits.

The selected improvement changes the normalization, which changes these
input monomials and allows removal of the stored $e$; it does not claim
to beat that restricted lower bound. Scaling instead by $\delta^2D$
would make the weights grow at real infinity and generally defeats the
preparation-zero bound $B<17$ without a further modification. That
variant is not used. The selected $\delta D$ scaling retains decay and
has been checked against the actual consuming inequalities above.
